\documentclass[12pt,a4paper]{article}
\usepackage[utf8]{inputenc}
\usepackage[T1]{fontenc}
\usepackage[brazil]{babel}
\usepackage{amsmath}
\usepackage{amssymb}
\usepackage{geometry}
\usepackage{hyperref}
\geometry{margin=2.5cm}

\title{Descricao Tecnica do Algoritmo de Reconhecimento de Quedas}
\author{Projeto Reconhecimento de Quedas}
\date{14 de abril de 2026}

\begin{document}
\maketitle

\section{Objetivo}
Este documento descreve o funcionamento do algoritmo de ponta a ponta, explicando:
\begin{itemize}
    \item como o sistema processa as entradas dos sensores;
    \item como decide entre \textbf{queda}, \textbf{instabilidade corporal} e \textbf{status normal};
    \item como cada sinal da camera contribui para a saida final de reconhecimento.
\end{itemize}

\section{Arquitetura do Sistema}
A aplicacao esta organizada em modulos no diretorio \texttt{src}:
\begin{itemize}
    \item \texttt{src/fall\_detection.py}: orquestracao principal;
    \item \texttt{src/fall\_core/camera.py}: captura de video e profundidade;
    \item \texttt{src/fall\_core/vision.py}: extracao de metricas de depth e esqueleto aproximado;
    \item \texttt{src/fall\_core/processing.py}: regras de deteccao de queda e risco;
    \item \texttt{src/fall\_core/state.py}: estado temporal;
    \item \texttt{src/fall\_core/events.py}: registro em CSV, snapshot e alerta.
\end{itemize}

\section{Entradas dos Sensores}
\subsection{Sensor RGB (color)}
Fornece o frame de video em formato BGR para:
\begin{itemize}
    \item visualizacao e sobreposicao de HUD;
    \item fallback de deteccao por pose 2D (MediaPipe), quando selecionado.
\end{itemize}

\subsection{Sensor de Profundidade (depth)}
Fornece um mapa de profundidade por pixel ($z$), usado como entrada principal no modo depth/skeleton para:
\begin{itemize}
    \item segmentar o individuo no primeiro plano;
    \item estimar caixa delimitadora do corpo;
    \item calcular centro corporal, altura relativa e razao de aspecto;
    \item estimar esqueleto simplificado (cabeca, quadril, pes) e angulo corporal.
\end{itemize}

\subsection{IMU da D435i (acelerometro e giroscopio)}
A D435i possui IMU, porem no estado atual do codigo a IMU \textbf{nao e usada diretamente} na inferencia.
Ela pode ser incorporada em evolucoes futuras para compensacao de vibracao da camera e fusao sensorial.

\subsection{Quais sao os 4 modulos frontais do dispositivo}
No uso pratico, muitos usuarios se referem a D435i como tendo ``4 cameras'' por causa dos 4 elementos opticos frontais visiveis.
Tecnicamente, eles sao:
\begin{itemize}
    \item \textbf{Camera IR esquerda}: participa da estimativa estereoscopica de profundidade;
    \item \textbf{Camera IR direita}: par estereo com a IR esquerda para calcular disparidade;
    \item \textbf{Camera RGB}: gera imagem colorida para visualizacao e algoritmos baseados em cor/pose 2D;
    \item \textbf{Projetor IR}: emite padrao infravermelho para melhorar textura da cena em superficies pobres de detalhe.
\end{itemize}

Importante: o projetor IR nao e uma camera de captura, mas faz parte do conjunto optico que melhora a qualidade do depth.

\section{Pipeline de Processamento (fim a fim)}
\subsection{1) Captura e alinhamento}
No modo RealSense, o pipeline captura streams RGB e depth e alinha depth para o frame de cor.

\subsection{2) Segmentacao do alvo humano}
A segmentacao e baseada em faixa de profundidade valida e operacoes morfologicas.
Em seguida, os contornos sao avaliados e o melhor alvo e escolhido por um escore area/profundidade.

\subsection{3) Extracao de metricas geometricas}
Do alvo principal, sao extraidas:
\begin{itemize}
    \item centro normalizado $(c_x, c_y)$;
    \item altura relativa $h_r$ e largura relativa $w_r$;
    \item razao de aspecto $AR = \frac{\text{largura}}{\text{altura}}$;
    \item profundidade mediana do torso.
\end{itemize}

\subsection{4) Estimacao de esqueleto aproximado}
A partir da mascara depth, o algoritmo estima:
\begin{itemize}
    \item cabeca: ponto mais superior;
    \item pes: ponto mais inferior;
    \item quadril: faixa central do corpo.
\end{itemize}
Com isso, calcula o angulo corporal em relacao a vertical:
\[
\theta = \arctan\left(\frac{|\Delta x|}{|\Delta y| + \epsilon}\right)
\]
quanto maior $\theta$, mais horizontal esta a postura.

\subsection{5) Suavizacao temporal}
Para reduzir ruido, o sistema aplica media movel exponencial (EMA):
\[
\hat{x}_t = \alpha x_t + (1-\alpha)\hat{x}_{t-1}
\]
com $\alpha = 0.35$ para sinais como centro de quadril, razao de aspecto e altura relativa.

\subsection{6) Classificacao do estado}
Por frame, o algoritmo gera um indicador de queda e um indicador de instabilidade.
A decisao final considera persistencia temporal e cooldown de eventos.

\section{Deteccao de Queda}
A logica de queda exige combinacao de postura, posicao no quadro e dinamica temporal.

\subsection{Condicoes principais}
\begin{itemize}
    \item \textbf{Postura deitada}: $AR$ alto ou angulo corporal elevado;
    \item \textbf{Centro baixo}: quadril mais proximo do chao na imagem;
    \item \textbf{Transicao de em pe}: historico recente de postura ereta;
    \item \textbf{Evento dinamico}: queda de altura relativa ou velocidade vertical brusca.
\end{itemize}

Formalmente, por frame:
\[
\text{fall\_frame} = \text{lying}\ \land\ \text{low\_center}\ \land\ \text{transition\_from\_upright}\ \land\ (\text{height\_dropped}\ \lor\ \text{sudden\_drop})
\]

\subsection{Confirmacao temporal de queda}
A queda so e confirmada apos acumulo de frames consecutivos de evidencias.
Apos confirmacao:
\begin{itemize}
    \item registra no CSV;
    \item salva snapshot;
    \item dispara alerta (se habilitado);
    \item aplica cooldown para evitar eventos duplicados imediatos.
\end{itemize}

\section{Deteccao de Instabilidade Corporal}
Instabilidade nao e queda imediata; e um estado de risco observado por oscilacoes.

\subsection{Metricas de risco}
No modo depth/skeleton, o algoritmo calcula:
\begin{itemize}
    \item \textbf{sway lateral}: amplitude em $x$;
    \item \textbf{velocidade lateral}: media de $|\Delta x|$;
    \item \textbf{jitter de quadril}: desvio padrao do centro vertical do quadril;
    \item \textbf{jitter angular}: desvio padrao do angulo corporal.
\end{itemize}

\subsection{Votacao multi-criterio}
Cada metrica acima do limiar gera 1 voto. O frame e marcado como instavel quando ha pelo menos 2 votos e nao ha queda no mesmo instante.

Depois, uma janela temporal de confirmacao exige repeticao em varios frames:
\[
\text{risk\_detected} = (\text{risk\_frames} \geq \text{risk\_confirm\_frames})
\]

\section{Reconhecimento de Status Normal}
O status normal e emitido quando:
\begin{itemize}
    \item nao existe queda confirmada;
    \item nao existe risco de instabilidade confirmado.
\end{itemize}

Logo, em termos logicos:
\[
\text{status\_normal} = \neg\text{fall\_detected} \land \neg\text{risk\_detected}
\]

\section{Mapa Entrada $\rightarrow$ Saida (consolidado)}
Esta secao resume, sem repeticao, como cada entrada participa da decisao final:
\begin{itemize}
    \item \textbf{IR esquerda + IR direita}: geram a base de profundidade estereoscopica que alimenta segmentacao corporal, postura e dinamica vertical;
    \item \textbf{Projetor IR}: melhora a confiabilidade do depth em cenas com baixa textura;
    \item \textbf{RGB}: usado para visualizacao/HUD e fallback por pose 2D;
    \item \textbf{IMU}: disponivel no hardware, ainda nao integrada ao classificador atual.
\end{itemize}

A partir dessas entradas, o sistema produz uma unica saida semantica por instante:
\begin{itemize}
    \item \textbf{Queda} (evento critico);
    \item \textbf{Instabilidade} (risco sem queda confirmada);
    \item \textbf{Normal} (ausencia de risco e queda).
\end{itemize}

\section{Registros e Evidencias}
A cada queda confirmada, o sistema gera:
\begin{itemize}
    \item linha no arquivo \texttt{relatorio\_quedas.csv};
    \item imagem em \texttt{capturas\_quedas/};
    \item tentativa de notificacao externa por script.
\end{itemize}

\section{Conclusao}
O algoritmo atual combina sinais geometricos de profundidade com memoria temporal para reconhecer estados de seguranca do individuo. O uso de depth torna a inferencia mais robusta que abordagens puramente RGB para cenarios de queda. A arquitetura modular permite evolucao para fusao com IMU/radar sem alterar o fluxo principal de decisao.

\end{document}
